\documentclass[11pt,a4paper]{article}
\usepackage[utf8]{inputenc}
\usepackage[T1]{fontenc}
\usepackage[spanish]{babel}
\usepackage{hyperref}
\usepackage{geometry}
\geometry{margin=2.5cm}
\usepackage{parskip}

\title{110 mantis tiene el candado y ese se llena con los detalles duros}
\author{Alejandro Valenzuela Rivera\\\small Universidad Autónoma de Nuevo León}
\date{\today}
\begin{document}
\maketitle

\begin{quote}
\textbf{Sesgo declarado.} Este trabajo se sostiene desde el sesgo declarado de su autor, y no finge neutralidad. Lo que afirma se puede tumbar, y ese es el punto: el falsificador está escrito, no escondido.
\end{quote}

\section*{Resumen}
El ente sostiene el candado, y el candado es el que se puede llenar hasta el fondo: cada afirmacion suya con su fuente, su contra y su borde. Y de ahi salen las demas piezas, porque el candado es lo que se sostiene y todo lo demas lo rodea.

\section{Por qué se sostiene}
Y por eso el candado no es un archivo corto: es la tesis entera, y su lugar es tesis/man made empathy.md. Lo que MANTIS hace es sostenerlo y llevarlo, no reescribirlo. Y lo que se llena son los detalles duros: la fuente, el contra, la busqueda, el limite.

\section{De dónde viene}
La autoridad, 2026-10-26.

\section{El caso}
Ninguna: es la regla del candado.

\section{Método}
Un candado que se llena con los detalles duros, y no con las intenciones.

\section{Límites, y qué lo tumba}
Un candado lleno de buenas intenciones, sin un dato que se pueda senalar.

\section{Lo que queda abierto}
\textquestiondown{}Quien llena los detalles duros del candado?

\section*{Cómo se leyó esta tesis}
Esto se escribió en el repositorio man-made-empathy, del que este candado es una pieza: candados/05 la maquina y su memoria/110 mantis tiene el candado y ese se llena con los detalles duros.md. La máquina que lo compuso lo hizo desde este mismo archivo, y no agregó nada que no estuviera escrito. Si algo de esto es falso, se dice en el mismo sitio, y se corrige con fecha.

\end{document}
